\documentclass{article}
\usepackage{amsmath,amssymb,booktabs}
\title{Disproof of Conjecture 00000001050:\\
A Relative-Trace Surjectivity Criterion with Nontrivial Kernel}
\author{TLMC submission}
\date{October 2, 2026}
\begin{document}
\maketitle

\section{The conjecture (verbatim)}

\begin{quote}
``Conjecture: For any composition of a permutation polynomial $f$ with a linearized
polynomial, the image under the relative trace $F_{q^2}/F_q$ is everything if and
only if the kernel is trivial (a relative-trace surjectivity criterion).''
\end{quote}

\section{Verdict}

\textbf{FALSE.} The biconditional fails already at $q = 2$ with the identity
permutation polynomial: the relative-trace image is everything while the kernel is
nontrivial.

\section{Counterexample}

Let $q = 2$ and $\mathbb{F}_4 = \mathbb{F}_2(\omega)$ with
$\omega^2 = \omega + 1$, $\omega^3 = 1$. Take
\[
f = \mathrm{id} \in \mathbb{F}_4[x] \quad\text{(a permutation polynomial)},\qquad
L(x) = \omega x + \omega x^2 = \omega(x + x^2) \in \mathbb{F}_4[x],
\]
where $L$ is a linearized polynomial (exponents $2^0, 2^1$; $\mathbb{F}_2$-linear
map). The composition is $g = L \circ \mathrm{id} = L$; since $f = \mathrm{id}$,
both composition orders coincide. The relative trace is
$\operatorname{Tr}_{\mathbb{F}_4/\mathbb{F}_2}(z) = z + z^2$.

\subsection{Pointwise computation}

\begin{center}
\begin{tabular}{cccc}
\toprule
$x$ & $x + x^2$ & $L(x) = \omega(x+x^2)$ & $\operatorname{Tr}(L(x)) = L(x) + L(x)^2$ \\
\midrule
$0$ & $0$ & $0$ & $0$ \\
$1$ & $1+1 = 0$ & $0$ & $0$ \\
$\omega$ & $\omega + \omega^2 = 1$ & $\omega$ & $\omega + \omega^2 = 1$ \\
$\omega^2$ & $\omega^2 + \omega^4 = 1$ & $\omega$ & $\omega + \omega^2 = 1$ \\
\bottomrule
\end{tabular}
\end{center}

\subsection{The two halves of the biconditional}

\begin{itemize}
\item \textbf{Kernel nontrivial.} $\ker L = \{x \in \mathbb{F}_4 : x + x^2 = 0\}
= \{0, 1\} = \mathbb{F}_2 \neq \{0\}$. (Indeed $x^2 = x$ has exactly the roots
$0, 1$ in any field.)
\item \textbf{Trace image is everything.} $L(\mathbb{F}_4) = \{0, \omega\}$, so
\[
\operatorname{Tr}_{\mathbb{F}_4/\mathbb{F}_2}\bigl(L(\mathbb{F}_4)\bigr)
= \{\operatorname{Tr}(0), \operatorname{Tr}(\omega)\}
= \{0, \omega + \omega^2\} = \{0, 1\} = \mathbb{F}_2 .
\]
\end{itemize}

Hence ``image is everything'' holds while ``kernel is trivial'' fails, so the
stated criterion
$\bigl(\operatorname{Tr}_{\mathbb{F}_{q^2}/\mathbb{F}_q}(g(\mathbb{F}_{q^2}))
= \mathbb{F}_q\bigr) \iff (\ker g = \{0\})$ is false. Conjecture 00000001050 is
disproven.

\section{Lean verification}

The file \texttt{lean4/Main.lean} encodes $\mathbb{F}_4$ as a four-constructor
inductive type with explicit addition/multiplication/Frobenius tables and proves:
\texttt{surjective\_true} ($\operatorname{Tr} \circ L$ hits $0$ and $1$),
\texttt{kernel\_not\_trivial} ($1 \in \ker L$, $1 \neq 0$), and
\texttt{criterion\_false} ($\neg(\text{Surjective} \leftrightarrow
\text{KernelTrivial})$). Every assertion is closed by \texttt{rfl} or constructor
exhaustion; there are zero axioms and no \texttt{sorry}. \texttt{lean4/Check.lean}
prints \texttt{\#print axioms} for every theorem.
An independent recomputation is available in \texttt{reproduce.py}.

\section{Boundary}

This counterexample only refutes the verbatim biconditional at the smallest
instance ($q=2$, $f = \mathrm{id}$). No claim is made about the converse direction
(kernel trivial $\Rightarrow$ surjective) in general, nor about compositions with
$f \neq \mathrm{id}$; both composition orders coincide here because $f=\mathrm{id}$.

\end{document}
